\chapter{Everyday operation}\label{everyday-operation}
\index{Shutdown}
\index{Command contexts}
\index{File retention}

Use this section as a route through the detailed, checked examples. Choose one fresh-session example at a time: do not create another process with a name already in use or enter FP commands while a demonstration program is waiting for an answer.

\section{Find the right procedure}\label{find-the-right-procedure}

\begin{longtable}[]{@{}
  >{\raggedright\arraybackslash}p{(\linewidth - 2\tabcolsep) * \real{0.5000}}
  >{\raggedright\arraybackslash}p{(\linewidth - 2\tabcolsep) * \real{0.5000}}@{}}
\toprule\noalign{}
\begin{minipage}[b]{\linewidth}\raggedright
Task
\end{minipage} & \begin{minipage}[b]{\linewidth}\raggedright
Procedure
\end{minipage} \\
\midrule\noalign{}
\endhead
\bottomrule\noalign{}
\endlastfoot
Start, enter the date and reach \passthrough{\lstinline!s!} & \manualref{first-monitor-session}{First monitor session} \\
Create a user FP, stop or remove a process & \manualref{processes-and-command-contexts}{Processes and command contexts} \\
Import, display, copy, rename or export text & \manualref{files-catalog-entries-and-algol-programs}{Files and Algol examples} \\
Edit an imported text file & \manualref{edit-a-text-file-inside-the-guest}{Guest editor} \\
Compile Algol or assemble SLANG & \manualref{compile-and-run-a-small-algol-program}{Algol example}, \manualref{assemble-a-slang-source}{SLANG example} \\
Find programs or inspect file metadata & \manualref{inspect-catalog-information}{Catalog information}, \manualref{software-and-files-in-the-supplied-disk-image}{installed software} \\
Preserve guest files for another session & \manualref{keep-guest-files-across-sessions}{Keep guest files} \\
Connect another terminal & \manualref{telnet-consoles}{Telnet consoles} \\
Choose text or binary host transfer & \manualref{choose-text-or-binary-transfer}{Transfer modes} \\
Run a reviewed FP command file & \manualref{call-one-command-file-from-another}{Command-file input}, \manualref{execute-commands-conditionally}{conditional execution} \\
Recover from accidental pasted input & \manualref{inspect-or-discard-pasted-input}{Inspect or discard input} \\
Select disk persistence or another controller & \manualref{choose-how-rc433x-disk-writes-survive}{RC433X persistence}, \manualref{select-the-optional-historical-disk-interface}{historical disk profile} \\
Limit host CPU use & \manualref{execution-throttle}{Execution throttle} \\
Inspect guest memory and execution & \manualref{debugging-and-core-files}{Debugging and core files} \\
\end{longtable}

\section{Decide what to keep before starting}\label{decide-what-to-keep-before-starting}

The default disk mode is disposable. Exporting a result with \passthrough{\lstinline!tohost!} creates a host file that survives emulator exit, but a guest file written only to the volatile disk overlay does not survive. For ordinary experiments, use that default and export wanted results before quitting.

For work that must remain on the guest disk, make both decisions explicitly:

\begin{itemize}
\tightlist
\item
  Select a persistent overlay or a writable working copy before booting.
\item
  Give the guest file a retained catalog scope using the checked \manualref{allocate-resources-and-promote-a-file}{resource allocation and promotion example}.
\end{itemize}

Persistent disk writes alone do not prevent the monitor from deleting a process's temporary files. Retained guest scope alone does not make a volatile disk overlay persistent. Follow the complete retention example rather than changing just one of these settings.

\section{Suspend, finish or exit}\label{suspend-finish-or-exit}

These are different actions, not interchangeable shutdown commands:

\begin{longtable}[]{@{}
  >{\raggedright\arraybackslash}p{(\linewidth - 4\tabcolsep) * \real{0.3333}}
  >{\raggedright\arraybackslash}p{(\linewidth - 4\tabcolsep) * \real{0.3333}}
  >{\raggedright\arraybackslash}p{(\linewidth - 4\tabcolsep) * \real{0.3333}}@{}}
\toprule\noalign{}
\begin{minipage}[b]{\linewidth}\raggedright
Action
\end{minipage} & \begin{minipage}[b]{\linewidth}\raggedright
Where
\end{minipage} & \begin{minipage}[b]{\linewidth}\raggedright
What remains
\end{minipage} \\
\midrule\noalign{}
\endhead
\bottomrule\noalign{}
\endlastfoot
Control-E, then \passthrough{\lstinline!continue!} & Simulator & Guest execution resumes without reset. \\
\passthrough{\lstinline!finis!} & Idle ordinary FP command input & Ends that FP; the system removes its process and temporary entries. \\
\passthrough{\lstinline!quit!} & \passthrough{\lstinline!sim4000>!} & Exits the emulator; volatile disk changes and running guest state are lost. \\
\end{longtable}

For a finished ordinary FP session, export wanted files first, enter \passthrough{\lstinline!finis!} and wait for the system's pause notice. If you need to verify removal, use a completed system \passthrough{\lstinline!list!}; the pause notice can precede cleanup. Then press Control-E and enter \passthrough{\lstinline!quit!} at the simulator prompt. This is not a rule for exiting an editor, an interactive diagnostic or another running program; use its own documented exit procedure to get back to FP first.

To take a break without ending the guest session, use Control-E and leave the emulator running at \passthrough{\lstinline!sim4000>!}, then \passthrough{\lstinline!continue!} later. The session is still only in that host process; quitting does not create a resumable snapshot. A saved core file is not a complete guest-session checkpoint.

\section{When the terminal context is unclear}\label{when-the-terminal-context-is-unclear}

At \passthrough{\lstinline!sim4000>!}, enter only simulator commands. At guest \passthrough{\lstinline!att!}, select a process name. In \passthrough{\lstinline!s!}, use system commands such as \passthrough{\lstinline!list!} or \passthrough{\lstinline!new!}. In an idle user FP, use utility names and file commands. A program-specific question is another context again: answer that program, not the host shell or FP.

If an expected system reply is delayed by \passthrough{\lstinline!to <process>!}, do not send another system command into that process's pending read. See the bounded \manualref{a-queued-fp-read-can-delay-a-system-reply}{queued FP read example}; its empty-line procedure applies only to the checked idle-FP situation. For other failures, use \manualref{startup-and-terminal-recovery}{startup recovery} or the \manualref{recover-from-fp-and-assembly-errors}{specific FP and assembly error examples}.
